% Modelo Canvas y análisis FODA (documento final para la feria)
\subsection{Modelo Canvas}\label{sec:feria_canvas}

El Modelo Canvas resume en nueve bloques cómo el sistema crea y entrega valor a sus usuarios. Como el sistema es institucional, la DUBSS lo financia con su presupuesto y no se vende a los estudiantes; por eso el bloque de ingresos se plantea como ahorro y eficiencia. La Tabla~\ref{tab:feria_canvas} presenta el modelo.

\begin{table}[H]
\caption{Modelo Canvas del sistema ComedorUAGRM}
\label{tab:feria_canvas}
\centering
\footnotesize\linespread{1}\selectfont
\renewcommand{\arraystretch}{1.15}
\setlength{\tabcolsep}{3pt}
\begin{tabular}{|>{\raggedright\arraybackslash}p{2.95cm}|>{\raggedright\arraybackslash}p{2.95cm}|>{\raggedright\arraybackslash}p{3.1cm}|>{\raggedright\arraybackslash}p{2.95cm}|>{\raggedright\arraybackslash}p{2.95cm}|}
\hline
\multirow{2}{*}{\parbox{2.95cm}{\textbf{Socios clave}\\[2pt] DUBSS (cliente y dueña del proceso); FICCT--UAGRM; concesionaria del comedor; Trabajo Social}} &
\textbf{Actividades clave}\newline Desarrollo web y móvil; entrenamiento y calibración del modelo de IA; auditoría y adjudicación; mantenimiento &
\multirow{2}{*}{\parbox{3.1cm}{\textbf{Propuesta de valor}\\[2pt] Postulación 100\% en línea; puntaje de vulnerabilidad con IA para todas las solicitudes; auditoría solo de finalistas; transparencia y trazabilidad; indicadores BI}} &
\textbf{Relación con clientes}\newline Autoservicio; consulta del estado de la solicitud; notificaciones; soporte de la DUBSS &
\multirow{2}{*}{\parbox{2.95cm}{\textbf{Segmentos de clientes}\\[2pt] Estudiantes postulantes; personal de la DUBSS; directivos; personal del comedor}} \\
\cline{2-2}\cline{4-4}
& \textbf{Recursos clave}\newline Equipo de desarrollo; módulo de IA; servidores (Docker); base de datos y archivos & & \textbf{Canales}\newline App móvil; panel web; correo y notificaciones push & \\
\hline
\multicolumn{2}{|>{\raggedright\arraybackslash}p{6.0cm}|}{\textbf{Estructura de costos}\newline Desarrollo del sistema; infraestructura y hospedaje; capacitación del personal; mantenimiento. Montos por definir con la DUBSS.} &
\multicolumn{3}{>{\raggedright\arraybackslash}p{9.0cm}|}{\textbf{Fuentes de ingreso / beneficio}\newline Financiamiento con el presupuesto institucional (no hay cobro al estudiante). Beneficio: ahorro de papel y de horas de revisión, y asignación más justa de unas 900 becas.} \\
\hline
\end{tabular}

\smallskip
\noindent\raggedright\textit{Nota.} Elaboración propia.
\end{table}

\subsection{Análisis FODA}\label{sec:feria_foda}

El análisis FODA contrasta los factores internos del proyecto (fortalezas y debilidades) con los externos (oportunidades y amenazas). La Tabla~\ref{tab:feria_foda} lo resume.

\begin{table}[H]
\caption{Análisis FODA del sistema ComedorUAGRM}
\label{tab:feria_foda}
\centering
\footnotesize\linespread{1}\selectfont
\renewcommand{\arraystretch}{1.2}
\begin{tabular}{|>{\raggedright\arraybackslash}p{7.4cm}|>{\raggedright\arraybackslash}p{7.4cm}|}
\hline
\textbf{Fortalezas (interno)} & \textbf{Debilidades (interno)} \\
\hline
\begin{itemize}\setlength\itemsep{0pt}\setlength\leftmargin{0pt}
\item Califica el 100\% de las solicitudes, no solo el 40\% que hoy se revisa.
\item Arquitectura modular en seis subsistemas.
\item Tecnologías abiertas y documentadas.
\item Trazabilidad con bitácora de acciones.
\item Apoyo directo de la DUBSS como cliente.
\end{itemize} &
\begin{itemize}\setlength\itemsep{0pt}\setlength\leftmargin{0pt}
\item El modelo de IA aún necesita datos históricos reales para calibrarse.
\item Puntaje mínimo de adjudicación y número de cupos por definir.
\item Tiempo de desarrollo limitado (24-ago a 30-nov de 2026).
\item Sin mecanismo de apelación en el alcance actual.
\end{itemize} \\
\hline
\textbf{Oportunidades (externo)} & \textbf{Amenazas (externo)} \\
\hline
\begin{itemize}\setlength\itemsep{0pt}\setlength\leftmargin{0pt}
\item Digitalización de los servicios universitarios.
\item Uso extendido del celular entre los estudiantes.
\item Posibilidad de integrarse con otros sistemas académicos de la UAGRM.
\item Replicable en otras becas o programas de bienestar.
\end{itemize} &
\begin{itemize}\setlength\itemsep{0pt}\setlength\leftmargin{0pt}
\item Resistencia al cambio del personal acostumbrado al proceso físico.
\item Estudiantes con acceso limitado a internet o a equipos.
\item Sensibilidad de los datos socioeconómicos (riesgo de filtración).
\item Cambios de autoridades o de presupuesto de la DUBSS.
\end{itemize} \\
\hline
\end{tabular}

\smallskip
\noindent\raggedright\textit{Nota.} Elaboración propia.
\end{table}
